\section{Discussion}\label{sec:disc}

\subsection{What transfers between earthquakes, and what does not}
Two signals are mixed in a coseismic landslide map: the \emph{relative} susceptibility of hillslopes to a given shaking field, and the \emph{absolute} level of landsliding in that event. Our results show that the first transfers reasonably well across very different Pacific Rim settings, and the second does not. Transferred models from all three learned families rank cells with AUC of about 0.73--0.89 in most events, and their calibration slopes near one show that the shape of the probability response is also approximately preserved. The base rate, by contrast, is mispredicted over a range of about three orders of magnitude, from $E/O=0.02$ at Cinchona to $E/O=23$ at Denali. Equal shaking and equal topography therefore produce very different numbers of landslides in different earthquakes. Some of this variation is genuine geoscience; some is a property of the inventories and ShakeMaps.

The events with the largest offsets have known, process-level explanations that our inputs do not include. At Denali the transferred models overpredict by a factor of more than 20. \citet{gorum2014denali} explicitly asked why this M\textsubscript{w}~7.9 earthquake triggered so few landslides, and \citet{jibson2004} used the sparse landsliding to infer the character of the strong shaking. Their explanations, discussed in those studies, involve rupture and site characteristics that a PGA and PGV field does not capture. At Niigata-Chuetsu ($E/O=0.23$) the earthquake struck weak Neogene sediments shortly after heavy typhoon rainfall \citep{sassa2005}. Antecedent pore pressure is outside our dry Newmark formulation and outside the static predictors. At Cinchona ($E/O=0.02$) the models predict almost no landslides because the ShakeMap Atlas gives a median PGA of only 0.19~g at the landslide cells and a maximum of 0.27~g in the domain. Yet thousands of slides, many of which became mudflows \citep{alvarado2009}, occurred on deeply weathered volcanic slopes. Either the ShakeMap underestimates near-source shaking for this shallow M\textsubscript{w}~6.1 event, or the volcanic soils fail at unusually low accelerations; most likely both contribute. At Porgera ($E/O=8.8$) the mapped domain is the largest of all events (24\,306~km$^2$), and dense tropical forest and persistent cloud may have limited the detectability of small slides. Mapping completeness thus probably adds to whatever physical difference exists \citep{tanyas2017,tanyas2022png}. These cases show that an event-level offset is itself a geoscientific quantity. It integrates antecedent wetness, material strength, shaking error and inventory completeness, and the transferred model cannot know it in advance.

The physics-based Newmark--Weibull curve calibrated at Northridge \citep{jibson2000} reflects the same issue. Its ranking skill (median AUC 0.76) is comparable to that of the statistical models, which shows that $D_N$ computed from generic strength and ShakeMap PGA already carries most of the transferable signal. Its probabilities are far too high (median $E/O=8.3$) and too spread out (calibration slopes 0.1--0.4). That curve encodes the Northridge inventory and its material properties, and it carries them to every other event. Re-estimating the displacement--probability relation on eight events (Physics-logit) reduces the bias but cannot remove the event-to-event variation.

\subsection{Physical constraints as protection against unphysical extrapolation}
Monotonic constraints did not raise the average transfer skill much, but they removed the failures. The global LR and, to a lesser extent, the unconstrained GBM failed exactly where the target event lay outside the training distribution: the very weak ShakeMap shaking of Ays\'en and Cinchona. In the LR the failure is analytic. A slope--PGV interaction that is well supported between 25 and 55~cm\,s$^{-1}$ implies, when extrapolated to 4~cm\,s$^{-1}$, that steep slopes are safer than gentle ones. Gradient boosting avoids sign reversals outside the data range because trees extrapolate as constants. Within the range, however, an unconstrained ensemble can still learn non-monotone responses to shaking from inventory noise, and these do not transfer. The constraint costs little where the data already support a monotone response (equivalent AUC at Kaik\={o}ura, Cinchona and Porgera). It loses at Denali and Palu, the two events where the observed landslide pattern is least consistent with ``more shaking, more landslides''. At Palu, many of the largest failures were long-runout flowslides on gentle irrigated slopes \citep{bradley2019palu}, a mechanism that a monotone slope response cannot represent.

The constraints follow the theory-guided approach of \citet{karpatne2017} and the call by \citet{reichstein2019} for physically consistent machine learning. Here they are encoded only as sign restrictions, a minimal form of prior knowledge that every geotechnical engineer would accept. More informative physics, such as event-specific strength or pore-pressure terms in the Newmark feature, could be added the same way.
